\section{Phase conventions and locality of the comparison fields}
\label{phases:section}

The gauges, links, and auxiliary comparison networks will all be realized
on one probability space.  We first specify how translations act on their
private inputs, then choose the wall phases and the comparison meshes.
These conventions give the column independence needed for the output
road and the finite-template reduction needed for two-dimensional
selected-path estimates.

\paragraph{Spatial private inputs and wall phases.}
Index a $d$-road by its structural role and its nominal origin modulo
$\Z d$, and index each private input also by its spatial location.
A unit longitudinal shift keeps the road-position class and shifts the
spatial index; a transverse or fractional-phase shift acts on both.
Auxiliary world starts and transition slots carry their own position
indices and translate as well. Distinct private coordinates are
independent, and all are independent of the edge weights. Their product
law is invariant under these shifts, as required by
Definition~\ref{def:fb-road}. Thus translated road copies
can use the same physical weights and independent private completions;
two distant portions of one road also use disjoint private coordinates
when their prescribed input neighborhoods are disjoint.

Fix a rational row spacing $\sigma$ such that every row increment in use,
including $s$, is an integer multiple of $\sigma$.  Let
$\mathcal H\subset\R/\Z$ be the finite subgroup generated by these
increments.  The canonical $B$-walls have nominal lines
$jA+\R B$, $j\in\Z$.  Independently of the spatial inputs, give these
walls independent uniform phases $\phi_j\in\mathcal H$, measured
relative to their nominal origins $jA$.  Their absolute phases in
$(v,B)$ coordinates are
\[
                         \theta_j=js+\phi_j\pmod 1.
\]
All auxiliary $A$-roads have deterministic longitudinal phases relative
to their nominal origins.  Their private randomness follows the spatial
indexing convention above.

Choose representatives $\phi_j\in[0,1)$.  A phase translates the road
rule together with its marked cuts and private inputs.  More explicitly,
if $\gamma_B(u)$ is the original $B$-road rule, the phased wall is
\[
 jA+\phi_jB+
 \gamma_B\bigl(\theta_{jA+\phi_jB}\omega,
                 \text{translated private inputs};u\bigr),
 \qquad t=u+\phi_j.
\]
Here $t$ is the nominal parameter relative to $jA$.  The marked cuts
$u=k$ have absolute heights $js+\phi_j+k$.  We restrict $n$ to
sufficiently divisible integers so that every phase displacement is a
lattice translation.  At each fixed phase, translation preserves the
product input law and the expected marked-period excess.  Consequently
the nonnegative period-defect estimates and all bounded-window
comparisons remain uniform after the phases are fixed.  Changing a phase
representative by an integer simply invokes the road's unit-period
covariance.

\paragraph{Averaging the comparison gauges.}
Fix one comparison radius $H$ large enough for every structural local
anchor attachment used below. In each offset mesh, register each labeled
anchor once at the lexicographically first nearest mesh index to its
deterministic nominal point, and reuse that registration wherever the
label appears. This rule commutes with translations and has a uniform
distance bound; longer comparisons telescope these local ones.
Lemma~\ref{found:lem-network} therefore permits this
choice before the canonical gauges are defined. Keep $H$, the cutoff
family, the reference anchors, and the offset group fixed when later
labels are added.

For each $a\in\mathcal H$, take the unit $(A,B)$ comparison mesh shifted
by $aB$, with auxiliary $A$-roads on its horizontal lines.  Apply
Lemma~\ref{found:lem-finite-cell} on this mesh. For each physical road
label used in a bounded comparison---a canonical $B$-wall, an $A$-road,
a corridor road, or a baseline introduced later---include the same
physical anchor as an additional label in every mesh. Average its gauge
functions over $a$. The formula~\eqref{found:eq-gauge-formula} for a
label uses only that anchor and the mesh's reference anchors, so these
labelwise definitions agree when two local comparisons overlap.

All labels taking part in a switch are averaged over the same offsets.
In particular, their averaged gauge difference is the average of the
within-mesh differences controlled by the local-envelope lemma. The
private inputs of all translated meshes are indexed as above, so
translations permute this finite family. This construction changes
only endpoint potentials: the physical roads and their anchor
occurrences are already fixed.

Adjoining a new road label in this fixed construction does not change
an existing gauge formula. Any new primitive coordinates placed in the
cells are jointly independent of the preexisting primitives;
conditioning an old variable on those extra coordinates
does not change its conditional expectation. Thus later baseline labels
can be added without changing the canonical wall gauges that define
the link field.

\begin{lemma}[Local phases and covariant gauges]\label{phases:locality}
The averaged gauges retain the local bounds, means, interpolation
identities, and common-tube comparisons of the preceding section.
Their support radii and bounds do not acquire a factor
$|\mathcal H|$.  They are covariant under translation by $A$, and their
laws are covariant under every row shift in $\mathcal H$.

For a link or gauge in a fixed neighborhood of wall $j$, all random phase
dependence is confined to a fixed number of wall indices near $j$.
Together with the independent spatial inputs, this gives independence
of entire columns at sufficiently separated indices.  The separation
constant is independent of $|\mathcal H|$ and $\sigma$.
\end{lemma}

\begin{proof}
We first check phase dependence, including that of the conditional
expectations used to construct a gauge.  An anchor is determined by
local weights, local private inputs, and the phases of the roads in a
fixed neighborhood.  The calibration itself depends on the weights
alone.  Each projected variable in
\eqref{found:eq-gauge-formula} is the difference of calibration values
at two anchors: the anchor carrying the gauge label and the reference
anchor indexed by $z_S$.  Both anchors lie within a fixed distance of
the gauge's location.  Indeed, for a gauge indexed by $z$, one has
$k(S)\in Q_R(z)$ and $z_S=k(S)+(R+4H+10)e_1$.

Regard the phases as fixed parameters when taking expectations over
weights and private inputs.  The variable being integrated depends on
the phases only through the two anchors.  Its conditional expectations
therefore depend on the same finite list of phases.  This applies to
every projection and to the deterministic anchor mean in the gauge
formula.  The number of potentially involved $B$-wall indices is
bounded by their spatial range, independently of the number of possible
phase values.  The spatial support of each projection is the one given
by Lemma~\ref{found:lem-finite-cell}.

All mesh offsets lie in one bounded fundamental cell.  Thus the supports
of the gauges being averaged lie in one fixed larger neighborhood,
uniformly in $|\mathcal H|$ and $\sigma$.  The triangle inequality for
the normalized average preserves the $L^1$ estimates without a factor
$|\mathcal H|$. At a fixed physical anchor, each mesh gauge has the
same mean, namely the expected calibration value at that anchor, so
averaging preserves this mean. For a switch between two labels, the
triangle inequality applied to the average of their within-mesh gauge
differences preserves the switch estimate without a factor
$|\mathcal H|$. Interpolation is affine in the anchor
gauges and uses the same physical road cost on every mesh.  It therefore
commutes with averaging.  The construction from the preserved incoming
road can be averaged in the same way; its common-tube comparison and its
dependence only on the baseline's own tube inputs and private field are
preserved.

Translation by $A$ advances the wall index by one and the absolute row
height by $s$.  Since $\theta_j=js+\phi_j$, its action on relative
phases is the ordinary shift of the sequence $(\phi_j)_{j\in\Z}$.
The road rules, spatial inputs, and gauge construction are covariant
under this joint translation, giving covariance under $A$.  Translating
the decorated network by $aB$, $a\in\mathcal H$, adds $a$ modulo one
to every wall phase and permutes the offset meshes.  Integer carries
are absorbed by unit-period covariance.  The independent uniform phase
law is invariant under this addition, and the spatial product law is
translation invariant.  This proves the asserted covariance of the row
laws.

Finally, the non-phase inputs for an entire column lie in a strip of
fixed width about that column.  Sufficiently separated columns have
disjoint strips and disjoint finite sets of wall-phase indices.  The
auxiliary $A$-roads contribute only private coordinates at the spatial
locations being inspected.  Hence the corresponding column
sigma-fields are generated by disjoint independent inputs.  The strip
width and the range of wall indices are independent of
$|\mathcal H|$ and $\sigma$, as required.
\end{proof}

\paragraph{Finite templates for selected-path estimates.}
A phase is shared along its entire wall, so the column independence
above does not give two-dimensional finite dependence. For each fixed
local test, however, Lemma~\ref{phases:locality} confines its phase
dependence to finitely many wall indices, including both anchors in
every gauge projection and the phases entering the anchor means.
A template is an assignment of values in $\mathcal H$ to this finite
list. Section~\ref{access:phases} extends these assignments periodically
on sparse grids. Each template then gives an iid field, and their
finite sum bounds the original nonnegative field pointwise. Different
template fields need not be independent.

All translated offsets used here are rational and are fixed before
$n$.  Sufficient divisibility of $n$ makes them lattice translations.
The number of templates may depend on the fixed auxiliary meshes and
on $s$, but it is independent of $n$ and of the eventual denominators
of the rational direction approximations.  This auxiliary enumeration
does not enter the input radius of the final road algorithm.
